\scanpage{67}
De cette façon, pour $\underline{\mathbb W}$ fixé, \deleted{$\unsure{\mathfrak X|\mathcal X}$ fixé} \deleted{$\mathbb W'$ \unsure{var}} $\mathfrak X$ et $\mathbb W'$ variables,
\[
 \mathbb W'{}^a(\mathfrak X\times_{\mathbb W(S)}\mathbb W'(S))
\]
devient un \uline{bifoncteur} en $\mathfrak X,\mathbb W'$, \uline{covariant} en $\mathfrak X$, \uline{contravariant en} $\mathbb W'$

\textcircled{2} \uline{Critère d'existence} : \uline{Cas affine}

% slanted
\margin{Ici, il suffit que $\mathbb W$ \unsure{au sens|au dessus} transforme affines en affines \deleted{\unsure{et soit fonc}}}
% slanted
\margin{Foncteur $\mathbb W$ \uline{spécial} [\unsure{associé} à un schéma \deleted{\unsure{d'}} \inserted{en} anneaux sur $S$ affine sur $S$]}

Supposons que $\underline{\mathbb W}$ transforme schémas affines (sur ouverts affines de $S$) en schémas affines. Donc, au dessus d'un ouvert affine de $S$, d'anneau $A$, $\underline{\mathbb W}$ est défini par un foncteur covariant, également noté $\underline{\mathbb W}$, transformant $A$-algèbres \inserted{$B$} en anneaux $\underline{\mathbb W}(B)$. [\unsure{Avec|Pour} l'hypothèse de \deleted{localité} \inserted{la nature locale} de ce foncteur, que nous n'\unsure{expliciterons|explicitons|explicitation} pas]. On suppose qu'on a $\underline{\mathbb W}(B)=\mathrm{Hom}_{A\text{-alg.}}(\Lambda,B)$ où $\Lambda$ est \deleted{l'algèbre affine} \inserted{l'anneau} d'un schéma en anneaux \deleted{affine} sur $S$, affine sur $S$.
% archive sheet number: 66
